\section{形式化结论及适用边界}\label{sec:supp-formal-zh}

本节给出当前证据能够支持的完整结论。证明刻意区分计账恒等式、容量松弛
下界、固定序证书与跨拓扑序全局最优性；后三者不能相互偷换。

\subsection{标准流量分解与两倍界}

对第 $i$ 次驻留中断，令 $v_i$ 为缓冲区大小；若该缓冲区无输入备份，令
$d_i=v_i$，否则令 $d_i=0$。定义
$V=\sum_i v_i$、$D=\sum_i d_i$。

\paragraph{命题 1（计账恒等式与紧界）。}
任一合法 P2/P3 工件的额外流量满足
\begin{equation}\label{eq:supp-identity-proof-zh}
 E=V+D,\qquad V\le E\le 2V.
\end{equation}
在固定 $V>0$ 时，仅改变被溢出字节的类别构成，两个工件的流量比至多为
2；该上界可达到。

\paragraph{证明。}
有备份缓冲区的一次中断收费 $v_i$，无备份缓冲区收费 $2v_i=v_i+d_i$。
逐项求和得
$E=\sum_i(v_i+d_i)=V+D$。又因每项 $d_i\in\{0,v_i\}$，有
$0\le D\le V$，故 $V\le E\le2V$。当所有被溢出字节均有备份时
$D=0,E=V$；当它们均无备份时 $D=V,E=2V$，两端均可达到。\hfill$\square$

命题 1 是评测器的计账性质，不是因果机制结论。它没有证明生产候选序读取
类别标签（事实上候选序生成不读取），也没有证明 $D$ 比 $V$ 更重要。任何
机理比较都应同时报告 $\Delta V$ 与 $\Delta D$；超过两倍的收益不可能在
固定 $V$ 下仅由类别构成产生。

\subsection{固定序驻留间隙模型}

固定一个合法原始拓扑序 $\pi$。对每个缓冲区 $b$，按 $\pi$ 中的位置排列
其强制事件：
$p_{b,0}<p_{b,1}<\cdots<p_{b,r_b}$。相邻事件之间的开区间定义一个驻留
间隙 $g=(b,j)$。二元变量 $x_g=1$ 表示 $b$ 从事件 $j$ 到事件 $j+1$
连续驻留，$x_g=0$ 表示允许中断该驻留。令
\begin{equation}
 w_g=\begin{cases}
 s_b,&b\text{ 有输入备份},\\
 2s_b,&b\text{ 无输入备份}.
 \end{cases}
\end{equation}

对存储类型 $k$ 和调度切面 $t$，令 $P_{k,t}$ 为该切面上强制必须驻留的
字节数，$\mathcal G_{k,t}$ 为跨过该切面的同类型可选间隙集合，容量为
$B_k$。固定序加权间隙模型为
\begin{align}
 L^\star=\min_x\quad &\sum_g w_g(1-x_g),\label{eq:supp-gap-obj-zh}\\
 \text{s.t.}\quad &P_{k,t}+\sum_{g\in\mathcal G_{k,t}}s_gx_g\le B_k,
 &&\forall k,t,\label{eq:supp-gap-cap-zh}\\
 &x_g\in\{0,1\}.\nonumber
\end{align}
若某个切面已有 $P_{k,t}>B_k$，则固定序本身不可由溢出修复，因为同一
强制事件所需的 pinned 集已经超过容量。

\paragraph{命题 2（固定序流量下界）。}
式~\eqref{eq:supp-gap-obj-zh}--\eqref{eq:supp-gap-cap-zh} 的最优值
$L^\star$ 不大于任意具有连续物理地址的合法固定序计划的流量。

\paragraph{证明。}
任取一个合法物理计划 $A$。若缓冲区在间隙 $g$ 内始终连续驻留，则令
$x_g=1$；否则令 $x_g=0$。在任一切面，$A$ 的强制驻留段与所有被选择的
连续驻留间隙确实同时占用该存储类型；因为 $A$ 合法，它们的总字节数不
超过 $B_k$，故构造出的 $x$ 满足所有容量约束。

每个 $x_g=0$ 的间隙至少包含一次完整驻留中断，因此 $A$ 在该间隙支付
至少 $w_g$；一个人为计划也可能在同一间隙内中断多次，支付更多。于是
$E(A)\ge\sum_gw_g(1-x_g)\ge L^\star$。该论证没有使用 $A$ 的具体地址，
故对所有合法连续地址计划成立。\hfill$\square$

\subsection{从下界到机器可检验证书}

容量松弛丢弃了地址区间的形状信息，所以即使允许从头全局安排所有地址，
其可行解也未必可打包。下面给出一个离散存储分配反例。容量为 4，七个段的
“活跃时间切片集合/大小”分别为
\[
 A:\{0\}/2,\ B:\{0,1\}/2,\ C:\{1,2\}/1,\
 D:\{1,2,3\}/1,\ E:\{2,3\}/1,\
 F:\{3,4\}/2,\ G:\{4\}/2.
\]
每个切片的活跃总量依次为 $4,4,3,4,4$，所以全部累计容量约束都成立；但
不存在一组全生命周期固定的连续 offset。证明如下：切片 0 强迫 $A,B$
各占一个长度 2 的半区间。不失一般性令 $B=[0,2)$。切片 1 又强迫 $C,D$
分别占 $[2,3)$ 与 $[3,4)$。若 $D=[2,3)$，则切片 2 的 $E$ 只能落在
$[0,2)$ 的某个单位位置；到切片 3，$D,E$ 把剩余空间割裂，长度 2 的 $F$
无处放置。若 $D=[3,4)$，只有 $E=[0,1)$ 时切片 3 尚可令
$F=[1,3)$；但到切片 4，$F$ 两侧各只剩 1，长度 2 的 $G$ 无处放置（其他
$E$ 位置已在切片 3 失败）。左右镜像同理。因此间隙模型只看总体积会接受
该组区间，“gap optimal”本身却不是物理可行或 canonical-valid 的证明。

\paragraph{定理 3（固定序流量证书）。}
对给定固定序，若同时满足：
\begin{enumerate}
  \item 间隙模型已经证明最优值为 $L^\star$；
  \item 所选间隙被实现为连续地址驻留段，且全部换出/换入依赖合法；
  \item canonical evaluator 返回零违规，并计算出 $E(A)=L^\star$；
\end{enumerate}
则工件 $A$ 在该固定序下具有最小额外流量。

\paragraph{证明。}
由命题 2，任一合法固定序物理计划 $A'$ 均满足 $E(A')\ge L^\star$；由条件
2--3，$A$ 是合法计划且 $E(A)=L^\star$。故不存在同一固定序下流量更小的
合法计划。\hfill$\square$

定理 3 的限定语不可省略。贪心打包或 \texttt{NoOverlap2D} 的
``OPTIMAL'' 只表示其可满足性模型证明找到/确认了布局，并非又优化了一次
流量；打包失败或超时不推翻 $L^\star$ 作为下界。该证书也不优化相同流量
下的溢出次数或时间，所以不能称为完整 P2 最优；固定序最优更不等于在所有
拓扑序上全局最优。

\subsection{候选组合的单调性}

令 $\mathcal A$ 为候选循环中成功生成的工件集合，$K$ 为实现内部按 P2 或
P3 语义计算的字典序键。生产求解器返回
$A^\star=\arg\min_{A\in\mathcal A}K(A)$。

\paragraph{命题 4（候选扩充的内部键单调性）。}
若新增成功生成的候选集合 $\mathcal B$，则
\begin{equation}
 \min_{A\in\mathcal A\cup\mathcal B}K(A)
 \ \le_{\mathrm{lex}}\
 \min_{A\in\mathcal A}K(A).
\end{equation}

\paragraph{证明。}
$\mathcal A\subseteq\mathcal A\cup\mathcal B$，右侧原有最优候选仍在新
集合中；在更大集合取最小值不可能更差。\hfill$\square$

这是 portfolio 对\emph{内部选择键}的单调性，而不是新候选有效性的因果
证据或合法性定理。要把返回值称为合法工件安全性，还需候选生成器本身正确，
或对最终返回工件运行 canonical evaluator；生产循环没有逐候选调用后者，
本文实验只报告最终验证为零违规的工件。特别地，若
组合包含一个强基线序，最终结果不劣于该序只是集合包含关系；仍需受控消融
才能把改进归因到某个排序、地址或 victim 组件。生产选择直接计算真实
$K_2/K_3$，不依赖 overflow-area 代理目标定理。

\subsection{复杂度与精确后端边界}

对固定候选数、victim 策略数和提前重载窗口数，可扩展路径依次执行拓扑
排序、驻留扫描、地址放置及时间模拟。当前放置器在一次请求上可能扫描驻留
集合或空闲区间，因此以缓冲区数计可保守写成二次项；portfolio 只乘一个
固定常数。该陈述是实现上界，不是线性时间保证。

加权间隙选择和后备 \texttt{NoOverlap2D} 在最坏情况下均为指数复杂度。
因此 CP-SAT 后端目前是离线研究 oracle，不属于 canonical
\texttt{solve}。将来若集成，必须设置规模阈值或超时，并在求解未完成、
打包失败或验证失败时回退到已经验证的可扩展工件。现有实验中的非零 gap、
timeout 和未运行项必须原样报告，不能用成功个案外推为普遍精确性。
